\chapter{The Region Frame: Reading the Model's Boxes Where It Drew Them}
\label{chap:s2-8}

\section{Goal}

Every region-based figure in Study~2 rests on reading the model's box correctly. While preparing the
model roster, a check of Qwen2.5-VL's documentation raised a question: does the model state its box in the
original photograph's pixels, or in the pixels of the resized image it was actually given? This chapter
records the answer, the fix, and every figure the fix changes.

\section{What Was Found}

Qwen2.5-VL resizes each image to fit the pixel budget (640 blocks of $28\times28$, about 0.5 megapixels)
before it sees it. A $1200\times900$ photograph becomes $812\times588$. The model states its box in that
resized frame. The harness read the box as original-image pixels, so every box was shrunk toward the
top-left corner by the resize factor, which was about 0.7 for most of the sample. Three independent checks
agreed:
\begin{itemize}
\item \emph{Edges.} On $1200\times900$ images the largest right edge the model ever wrote is 839, just
past the resized width of 812 and far short of 1200.
\item \emph{Overlap.} Mapping the boxes back to the original frame raises mean IoU with the annotated
boxes from 0.058 to 0.389 (4-bit) and from 0.060 to 0.398 (bf16). IoU $\geq 0.5$ goes from 2.1\% to
40.6\% (\texttt{results/diag\_region\_frame.txt}).
\item \emph{Processor.} Running the model's own image processor on a $1200\times900$ blank gives a
$29\times21$ patch grid, i.e.\ $812\times588$ pixels: exactly the frame the boxes fit.
\end{itemize}

\textbf{Why the earlier check missed it.} Before F3 was reported, the highest-IoU rows were checked by
hand against the annotation and matched to three decimals. That check was sound for the rows it looked at,
and those rows were the problem. Eighteen of the 25 rows with IoU $\geq 0.5$ were images already within
the pixel budget, where the resize factor is about 1 and the two frames coincide. Such images are about 5\%
of the sample. Checking only the rows that looked right selected exactly the rows the error did not
touch. A coordinate check needs rows chosen without regard to their score.

\section{What Was Done}

\textbf{The fix, behind a flag.} \texttt{xai\_vlm/frames.py} maps a box from the model-input frame to the
original image, using the patch grid the processor reports for that image (\texttt{image\_grid\_thw},
14-pixel patches). The backend applies it when \texttt{XAI\_REGION\_FRAME=rescaled}. The default,
\texttt{legacy}, reads boxes as before, so every frozen file still reproduces byte-for-byte. A rescaled
run records the flag in its stage fingerprints and must write to its own results tree. An empty value
counts as unset; the first validation job failed on exactly that and is kept in the log.

\textbf{Stages that only read regions: no GPU.} \texttt{17\_region\_frame\_fix.py} builds
\texttt{results/scholar-a40-rf/} from the frozen A40 tree without touching it:
\begin{itemize}
\item Stage~2 regions are re-parsed from each stored reply in the input frame, then mapped. Nothing is
lost or gained: 4-bit 3,022 mapped and 860 empty, bf16 3,086 and 796.
\item Stage~6 rerun regions and Stage~8 perturbed regions are mapped; drift is recomputed.
\item Stages~3, 7 and 10 then re-run against the new tree.
\end{itemize}

\textbf{Stages that used the regions on the GPU.} Stage~4 masks the region and Stage~5 measures attention
inside it, so both re-ran on the A40 for both tiers under the flag, followed by Stages~9, 10 and 11.

\textbf{Validation.} Before the reruns, \texttt{cluster/rf\_validate.sbatch} re-asked 40 rows with the
backend mapping boxes itself. All 40 replies matched the frozen replies word for word, and all 40 mapped
boxes matched the offline mapping (\texttt{RF\_VALIDATE PASS}, job 488128). Every new Stage~4 row uses
exactly the corrected Stage~2 region, and Gini, which does not depend on the region, is unchanged
(median 0.555 in both trees).

\textbf{Two small gaps.} Stage~10 re-ran on the GPU node: two 4-bit rows and one bf16 row differ from the
CPU build in the fourth decimal of text similarity, a numerical difference with no effect on any verdict.
Stage~5 was meant to save its attention grids so that later region changes could be re-scored without a
GPU. The cluster copy predated that change, so no grids were saved. The figures are unaffected.

\section{Results}

\begin{table}[tbp]
\centering
\small
\begin{tabular}{p{6.0cm}rrrr}
\toprule
 & \multicolumn{2}{c}{4-bit} & \multicolumn{2}{c}{bf16} \\
\textbf{metric} & \textbf{before} & \textbf{after} & \textbf{before} & \textbf{after} \\
\midrule
\multicolumn{5}{l}{\emph{Stage 10: is the explanation correct? (true YES rows)}} \\
explanation correct & 2.4\% & 31.6\% & 2.1\% & 33.8\% \\
IoU $\geq 0.30$ & 6.7\% & 54.7\% & 6.6\% & 55.8\% \\
median IoU & 0.00 & 0.37 & 0.00 & 0.39 \\
right answer, wrong reason & 97.6\% & 68.4\% & 97.9\% & 66.2\% \\
\multicolumn{5}{l}{\emph{Stage 4: does masking the region flip the answer?}} \\
flip, black mask (YES rows) & 6.4\% & 25.8\% & 6.8\% & 27.9\% \\
artefact floor, black (NO$\to$YES) & 27.2\% & 24.4\% & 31.6\% & 33.8\% \\
flip, token ablation (YES rows) & 6.1\% & 15.8\% & 5.2\% & 16.1\% \\
artefact floor, tokens (NO$\to$YES) & 51.1\% & 40.3\% & 58.2\% & 45.1\% \\
\multicolumn{5}{l}{\emph{Stage 5: is attention inside the region?}} \\
attention exceeds area null & 85.4\% & 97.2\% & 85.9\% & 97.6\% \\
median attention mass inside & 2.4\% & 23.8\% & 2.2\% & 23.6\% \\
median region area & 1.1\% & 2.4\% & 1.0\% & 2.1\% \\
\multicolumn{5}{l}{\emph{Stage 9: completeness, attention basis}} \\
explanation supported & 66.5\% & 75.7\% & 68.3\% & 77.6\% \\
lift $\geq 2\times$ (regions only) & 42.6\% & 77.5\% & 43.2\% & 78.7\% \\
\multicolumn{5}{l}{\emph{Stages 6 and 8: does the region stay put?}} \\
rerun drift $>10\%$ of diagonal & 36.1\% & 56.2\% & 38.2\% & 55.4\% \\
blur: verdict holds, region moves & 47.7\% & 61.2\% & 42.6\% & 57.4\% \\
occlusion: verdict holds, region moves & 35.1\% & 46.3\% & 34.2\% & 46.2\% \\
\bottomrule
\end{tabular}
\caption{Every region-based figure, A40 tree before (\texttt{scholar-a40}) and after
(\texttt{scholar-a40-rf}) the fix. Verdict-only figures (sensitivity, specificity, stability of the
verdict, survival under perturbation) do not change. Intervals are in
\texttt{results/scholar-a40-rf/report\_<tier>.txt}.}
\label{tab:s28}
\end{table}

\textbf{The model locates hazards far better than reported.} When it correctly says Yes, its box overlaps
the annotated box by IoU 0.3 or more 55\% of the time, and a third of its explanations are fully correct
(right hazard, similar reason, overlapping box). F3's headline of 2.4\% was the coordinate error, not the
model. The weak point is now split between the box and the wording: 56\% of reasons pass the similarity
bar and 55\% of boxes pass the IoU bar. Poor-info images score higher than rich-info ones (43\% against
24\% for 4-bit), a split that needs explaining.

\textbf{Attention sits inside the box.} Median attention mass inside the stated region is ten times the
region's share of the image (23.8\% against 2.4\%), and it beats the area null on 97\% of rows. On the
attention basis, 78\% of explanations have a lift of $2\times$ or more, against 43\% before. The old
tension between F3 and F6 (attention on target, boxes off target) disappears: the boxes were on target
and the harness misplaced them.

\textbf{Masking the region now moves the answer, but only some rules clear the artefact floor.} Blacking
out the corrected region flips a quarter of Yes answers (25.8\% [24.2, 27.5] for 4-bit), four times the
old figure. The black-mask floor is 24.4\% [20.4, 28.7], so overall the signal and the floor still
overlap. By rule they separate:
\begin{itemize}
\item rule\_4: 65.2\% [53.0, 75.8] (4-bit) and 54.6\% (bf16), well above the floor;
\item rule\_2: 45.5\% [40.2, 50.8] and 40.8\%, above it;
\item rule\_1: 30.6\% [28.0, 33.2] and 28.7\%, at the floor;
\item rule\_3: 9.8\% and 15.6\%, below it.
\end{itemize}
Token ablation still has a floor (40.3\%) far above its signal (15.8\%), so the token path remains
uninformative.

\textbf{Regions are less stable than reported.} Measured in the right frame, boxes move further. Between
greedy and sampled reruns, 56\% drift by more than a tenth of the image diagonal (36\% before). Under mild
blur the verdict often holds, but the box moves that far 61\% of the time.

\section{Standard vs.\ Non-Standard Choices}

\textbf{Standard.} Reading a model's coordinates in the frame it was given; Qwen2.5-VL's documentation
uses absolute coordinates in the resized input.

\textbf{Non-standard, and stated in the paper.}
\begin{itemize}
\item The fix is behind a flag and the frozen trees are kept, so every pre-fix number can still be
reproduced and compared. The paper reports the corrected tree.
\item The stages that only read regions were corrected offline from stored replies rather than re-run.
That is exact for Stage~2 (the reply is kept) and was confirmed against a fresh GPU run on 40 rows.
\end{itemize}

\section{Implications}

\begin{itemize}
\item \textbf{F3 reverses.} The model does not ``describe hazards well and locate them badly''. It
locates the annotated hazard about half the time when it is right, and its full explanation is correct
about a third of the time. ``Right answer, wrong reason'' falls from 98\% to 68\%: still the majority,
but no longer near-total.
\item \textbf{F4 softens and becomes rule-specific.} The probe is not drowned everywhere. For struck-by
(rule\_4) and fall protection (rule\_2), masking the stated region removes the answer well beyond what
masking does to answers that had nothing to find. For PPE it does not separate, and for edge protection
it is below the floor.
\item \textbf{F6 strengthens, and F7's lift caveat is gone.} Attention is inside the box on 97\% of rows
at a median lift of $8\times$ (4-bit) and $9\times$ (bf16); over three quarters of regions are at
$2\times$ or better.
\item \textbf{Region stability is a new weakness.} The box is right more often than reported, and less
steady than reported.
\item \textbf{For the model roster.} Every new model states boxes in its own convention (resized pixels,
0--1000, 0--1, points). Each adapter will map to the original frame and be checked on rows drawn without
regard to their score, before any metric is computed.
\end{itemize}
